\section{Conclusion}
%================================================
\begin{frame}{Conclusions}

	\begin{block}{Main Findings}
		\begin{itemize}
			\item Experimental droplet trajectories exhibit persistent self-avoiding motion.

			\item MSD analysis indicates superdiffusive transport
			      ($\alpha > 1$).

			\item VACF and OCF reveal finite translational
			      ($\tau_p$) and rotational ($\tau_r$) memory.

			\item A memory-driven PDE--SDE model was developed
			      based on droplet--chemical field interactions.

			\item The model reproduces the key qualitative
			      transport and memory signatures observed experimentally.
		\end{itemize}
	\end{block}

	\begin{alertblock}{Key Limitation}
		The model reproduces the qualitative decay of the VACF and OCF,
		but significantly underestimates the persistence times
		($\tau_p,\tau_r$) measured experimentally.
	\end{alertblock}

\end{frame}
%================================================
% \begin{frame}{Future Work}

% 	\begin{itemize}

% 		\item Calibrate model parameters using experimental measurements
% 		      to enable quantitative comparison.

% 		      \vspace{0.2cm}

% 		\item Investigate alternative memory kernels and generalized
% 		      Langevin descriptions.

% 		      \vspace{0.2cm}

% 		\item Study collective behaviour and interactions between multiple
% 		      active droplets.

% 		      \vspace{0.2cm}

% 		\item Extend the framework to other memory-driven active matter
% 		      systems.

% 		      \vspace{0.2cm}

% 		\item Explore long-time transport and anomalous diffusion regimes.

% 	\end{itemize}

% 	\vspace{0.5cm}

% 	\begin{alertblock}{Long-Term Goal}
% 		Develop predictive memory-driven models for active matter systems
% 		that couple stochastic motion with evolving environmental fields.
% 	\end{alertblock}

% \end{frame}

\begin{frame}{Outlook and Future Work}

	\begin{block}{Physical Interpretation}
		The results support the hypothesis that active droplet
		motion is influenced by memory effects arising from
		self-generated chemical fields, leading to
		history-dependent dynamics.
	\end{block}

	\begin{block}{Future Directions}
		\begin{itemize}
			\item Calibrate model parameters for quantitative agreement.

			\item Investigate memory-kernel based descriptions
			      (e.g., Generalized Langevin Equation).

			\item Extend the framework to interacting droplets
			      and collective behaviour.

			\item Explore long-time transport and anomalous
			      diffusion regimes.
		\end{itemize}
	\end{block}

	\begin{alertblock}{Long-Term Goal}
		Develop predictive memory-driven models for active matter
		systems that couple stochastic motion with evolving
		environmental fields.
	\end{alertblock}

\end{frame}
